\documentclass[11pt]{article}
\usepackage[T1]{fontenc}\usepackage{lmodern,amsmath,amssymb,booktabs}
\usepackage[margin=1in]{geometry}\usepackage[hidelinks]{hyperref}
\setlength{\parindent}{0pt}\setlength{\parskip}{5pt}
\newcommand{\E}{\mathbb E}\newcommand{\Var}{\operatorname{Var}}\newcommand{\KL}{\operatorname{KL}}
\begin{document}
\begin{center}{\LARGE\bfseries Marginal bias and exact sampling in\\[3pt] finite determinantal candidate pools}\\[7pt]
Lluis Eriksson\\[2pt]\small Working manuscript, 20 November 2026\end{center}
\textbf{Abstract.} A language model can generate a pool of candidates and then sample a diverse subset using a determinantal point process. For a fixed finite catalog we compute the first two possible leading terms of the divergence between the target distribution and the sampler's marginal output law. The generic $N^{-2}$ coefficient depends on first-order inclusion probabilities. Matching the proposal to these inclusions removes that coefficient, and the $N^{-4}$ coefficient then depends on pair inclusions. For two-item determinantal samples with redraw, either this second coefficient is positive or the sampler is exact at every pool size $N\ge2$. A three-item example gives coefficients $5/72$ before matching and $4/25$ after matching, at their respective orders. An explicit six-line example shows that the two-item dichotomy fails for three-item samples.

\section{The output of a random pool}
Prada-Corral et al.\ \cite{pool} study sampling a set of $K$ distinct strings. They draw a pool of $N$ strings independently from a proposal $q$ and sample from its importance-weighted $K$-DPP. Their Theorem 5.2 gives a finite-sample upper bound on forward KL, including for countable supports under their moment conditions. Here the catalog, $q$, $K$ and the target $P$ are fixed while $N\to\infty$. Constants may depend on the entire instance. The method of weighted resampling is established \cite{kviman}; the issue here is which inclusion probabilities determine its marginal bias.

Let $P$ be a probability law on $K$-subsets of a finite catalog. Retain only active items and assume that $q$ is positive on them, sums to one, and $N\ge K$. In the determinantal case $P(S)$ is proportional to the squared volume of $K$ feature vectors \cite{kt}. Let $n_i$ count occurrences of item $i$. Define
\[
 \delta_i=\frac{n_i}{Nq_i}-1,\qquad
 w_S=\prod_{i\in S}(1+\delta_i),\qquad Z=\sum_T P(T)w_T.
\]
Conditional on $Z>0$, output $S$ with probability $P(S)w_S/Z$. Each occurrence contributes a choice of copy, while duplicate columns selected together have zero determinant. This recovers the pool law in \cite{pool}. Redraw the pool when $Z=0$, and write $Q_N$ for the resulting \emph{marginal} distribution of one set. All logarithms are natural. For a particular pool, a positive target set can be missing entirely.

\newpage
\section{The inclusion coefficients}
Let $p_i=P(i\in S)$ and $p_{ij}=P(i,j\in S)$ for distinct $i,j$. Use the two scores
\[
 T_S=\sum_{i\in S}p_i/q_i,\qquad D_S=\sum_{i<j\in S}p_{ij}/(q_iq_j).
\]
\textbf{Theorem.} Uniformly over the fixed support of $P$,
\[
 \frac{Q_N(S)}{P(S)}=1+\frac{\E_P T-T_S}{N}+O(N^{-2}),\qquad
 \KL(P\|Q_N)=\frac{\Var_P(T)}{2N^2}+O(N^{-3}).
\]
For the matched proposal $q_i=p_i/K$,
\[
 \frac{Q_N(S)}{P(S)}=1+\frac{\E_P D-D_S}{N^2}+O(N^{-3}),\qquad
 \KL(P\|Q_N)=\frac{\Var_P(D)}{2N^4}+O(N^{-5}).
\]
These expansions use the multiplicity formula above; they do not require that $P$ itself be determinantal. For $K=1$ the empty pair sum is zero and matching gives the target exactly.

\textit{Expansion at the population frequencies.} Write $L_S=\sum_{i\in S}\delta_i$, $B_S=\sum_{i<j\in S}\delta_i\delta_j$, and let $V_j$ denote the degree-$j$ part of $Z$. Thus $V_1=\sum_i p_i\delta_i$ and $V_2=\E_P B_S$. Each $\delta$ is a mean of $N$ centered i.i.d.\ vectors $X$ with $X_i=\mathbf1_{\{I=i\}}/q_i-1$, $I\sim q$. Its one-draw covariance is $C_{ij}=\mathbf1_{\{i=j\}}/q_i-1$. Multinomial factorial moments give $\E w_S=(N)_K/N^K$ for every $S$, hence $\E(w_S-Z)=0$.

Expand $(w_S-Z)/Z$ near zero. The first surviving term in expectation at order $N^{-1}$ is $V_1^2-L_SV_1$. Its expectation times $N$ is $\E_P T-T_S$. For justification, all fixed coordinates concentrate exponentially in a fixed neighbourhood where $Z$ stays bounded away from zero. The bounded conditional probabilities and all polynomial terms contribute exponentially little off that event. A degenerate pool is exponentially rare, since one positive target set is present except with probability at most $\sum_{i\in S}(1-q_i)^N$. Redrawing has no effect on the displayed polynomial coefficients. Taylor expansion to degree three gives the first remainder: centered third joint moments are $O(N^{-2})$ and fourth absolute moments are $O(N^{-2})$.

\newpage
\textit{Matched proposal.} Now $V_1=K\sum q_i\delta_i=0$ identically. Up to terms of degree at least five, the correction to the numerator $w_S-Z$ in the ratio is
\[
       -L_S V_2 -(B_S-V_2)V_2 - L_S V_3 .
\]
Terms of degree above $K$ in $Z$ are understood to vanish. For a single draw, $V_2(X)=\binom K2-(K-1)p_I/q_I=-\binom K2$ is constant. The expected cubic term is $N^{-2}\E[L_S(X)V_2(X)]=0$ by centering and independence across draws.

The $N^{-2}$ coefficient of a fourth-degree moment is computed by the three covariance pairings. Terms where all four draws coincide contribute only $O(N^{-3})$. In $L_S V_3$, each term of $V_3$ has three distinct indices. Pair two of them: their covariance is $-1$. For a fixed index $h$ of $L_S$, summing the remaining covariances gives $-\binom{K-1}{2}\sum_i p_i C_{hi}=0$. This also covers $K=2$ where $V_3=0$.

To evaluate the middle product, put $A_{ij}=p_{ij}/2$ off the diagonal, $A_{ii}=0$. Then $V_2=\delta^\top A\delta$. The row sums are $(K-1)p_i/2=\binom K2 q_i$. Consequently $(CAC)_{ij}=p_{ij}/(2q_iq_j)-\binom K2$ for $i\ne j$. In the pairing computation for $B_S V_2$, the terms pairing each factor internally cancel against those for $V_2^2$, since $\E B_S=\E V_2$. The two cross-pairings sum to $2\sum_{i<j\in S}(CAC)_{ij}$ for $B_S V_2$, and to its $P$-average for $V_2^2$. Their difference is $\E_P D-D_S$ with the sign in the displayed correction.

Taylor expansion through degree five leaves a sixth absolute moment of order $N^{-3}$; a centered fifth joint moment also has order $N^{-3}$ since a nonzero index pattern has at most two distinct draw indices. This gives the claimed matched remainder. Finally set $r_S=Q_N(S)/P(S)-1$. Since $\E_P r_S=0$ with redraw, the linear term of $-\log(1+r_S)$ cancels exactly. Applying the expansion to the remaining quadratic and higher terms proves both KL formulas. \hfill$\square$

\textbf{Which proposals cancel the first coefficient?} For a $K$-DPP, connect items if exchanging them in a positive set can yield another positive set. Constancy of $T_S$ implies $p_i/q_i$ is constant along each edge. Conversely positive sets are connected by single exchanges: from a basis $A$ toward $B$, use a vector of $B$ outside the span of $A$, if available; otherwise a vector of $B$ uses a nonzero coordinate of $A\setminus B$ in that basis. Each swap can increase the intersection. Thus componentwise constancy suffices. For a connected exchange graph the unique canceling proposal is $p_i/K$.

\newpage
\section{Two-item samples: fourth order or exact}
Suppose $K=2$, $q_i=p_i/2$, and $P$ is a finite $2$-DPP. Delete inactive items. Partition the remaining nonzero feature vectors into classes of parallel directions, disregarding sign. A pair is positive exactly when its items belong to different classes.

\textbf{Corollary.} Under redraw, either $\Var_P(D)>0$ and forward KL is asymptotic to a positive constant times $N^{-4}$, or $Q_N=P$ for every $N\ge2$.

\textit{Proof.} If the variance vanishes, the pair score $D_{\{i,j\}}=P(\{i,j\})/(q_iq_j)$ is a common constant $c$ on the positive pairs. If $i$ is in class $g$ with proposal mass $m_g$, summing over its possible partners yields $2q_i=cq_i(1-m_g)$. All classes therefore have equal mass. Conditional on pool counts, the probability of choosing a class pair $g,h$ is proportional to $n_g n_h$. The vector of class counts is exchangeable, also after conditioning on having at least two occupied classes. Every unordered pair of classes has the same marginal selection probability. Given the class counts, within-class labels are independent multinomial samples from $q_i/m_g$; selecting a random occurrence in each chosen class thus gives precisely those conditional item probabilities. The resulting pair law is proportional to $q_iq_j$ across distinct classes, hence equals $P$. The positive-variance case is the theorem. \hfill$\square$

\section{A boundary at three-item samples}
Use six equiangular lines in $\mathbb R^3$. With $\phi=(1+\sqrt5)/2$, take the columns
\[
 (0,\pm1,\phi),\quad (\pm1,\phi,0),\quad(\phi,0,\pm1)
\]
and normalize each to length one. Direct multiplication gives $BB^\top=2I$ and every off-diagonal entry of $G=B^\top B$ has magnitude $1/\sqrt5$. Thus $\Pi=G/2$ is a rank-three orthogonal projection. Use its three-item projection DPP. By Cauchy--Binet, inclusion probabilities are the principal minors of $\Pi$: expand $\det(I+B\,\mathrm{diag}(t_i)B^\top/2)$ in the $t_i$, or expand $\det(B\,\mathrm{diag}(1+t_i)B^\top/2)$ in bases and compare coefficients. Hence $p_i=1/2$, $p_{ij}=1/5$, $q_i=1/6$ and $D_S$ is constant. Both displayed KL coefficients vanish.

Yet each triple has probability $(1\pm1/\sqrt5)/20$ according to the sign of its three pairwise inner products multiplied together. Both signs occur: with the columns ordered by taking the plus sign first in each pair, triples $(1,2,3)$ and $(1,3,5)$ have opposite signs. All triples are positive. At $N=K=3$ a successful pool consists of three distinct items, so its unordered triple is uniform over all twenty possibilities. This contradicts exactness. The theorem implies $Q_N/P-1=O(N^{-3})$ and hence $\KL(P\|Q_N)=O(N^{-6})$ for this example. A sharp sixth-order coefficient is not inferred from the small-pool discrepancy.

\newpage
\section{Checks and implications}
\textbf{Example.} For diagonal kernel entries $(6,3,2)$ the three pair probabilities in lexicographic order are $(1/2,1/3,1/6)$. Uniform proposals give leading KL coefficient $5/72$ at order $N^{-2}$. Matching gives $q=(5/12,1/3,1/4)$ and pair scores $(18/5,16/5,2)$. The relative bias coefficients at order $N^{-2}$ are $(-2/5,0,6/5)$. Thus the leading KL coefficient is $4/25$ at order $N^{-4}$. The companion checker enumerates every count vector through selected pool sizes up to 64, and computes the two Taylor coefficients independently by expanding the conditional law as a polynomial. A nonuniform parallel-class control checks exact sampling under redraw.

The \href{https://github.com/lluiseriksson/THE-ERIKSSON-PROGRAMME/tree/codex/pool-bias-20261120/audits/pool-dpp}{source and companion checker} also include a diagonal three-item-sample control. Exact rational arithmetic is used for the small-catalog count enumeration, with decimal logarithms only in the final divergence calculation.

\textbf{Scope.} The marginal inclusions may themselves be expensive to compute. This is a criterion for analyzing proposals on fixed catalogs. It gives neither a distribution-free choice of pool size nor a guarantee for greedy MAP, growing supports, or language-model task accuracy. No independent empirical evaluation on model outputs was performed.

\section*{Model assessments and review context}
OpenAI Codex assisted with derivation, writing and code in a shared authoring context. Assessment of the exact PDF is recorded separately. Independent peer review has not been documented.
\begin{thebibliography}{3}\setlength{\itemsep}{0pt}\small
\bibitem{pool} M. Prada-Corral et al. A Model of Diverse Sampling from Language Models. \href{https://arxiv.org/abs/2611.25003v1}{arXiv:2611.25003v1} (2026), Sections 4--5 and Appendix A.
\bibitem{kviman} O. Kviman et al. Variational Resampling. \href{https://proceedings.mlr.press/v238/kviman24a.html}{AISTATS}, 2024.
\bibitem{kt} A. Kulesza and B. Taskar. $k$-DPPs: Fixed-Size Determinantal Point Processes. \href{https://icml.cc/2011/papers/611_icmlpaper.pdf}{ICML}, 2011.
\end{thebibliography}
\end{document}
